%
\documentclass[11pt]{article}

\usepackage{times}
\usepackage{fullpage}
\usepackage{amsmath,amssymb,bm}
\usepackage{booktabs}
\usepackage{tabularx}
\usepackage{hyperref}
\usepackage{enumitem}
\usepackage{setspace}

\title{An LLM-Run Quarterly Firm Laboratory: Bankruptcy/Entry, Compustat-Style Data, and Debrief Scoring (Draft)}
\author{}
\date{}

\begin{document}

\maketitle

\begin{abstract}
This document specifies a repeated-game ``laboratory'' where firms are controlled by language models (LLM agents) and interact through (i) a \emph{product market} and (ii) a \emph{financial market} (an LLM agent). The design goal is (a) a realistic-enough operating, financing, and accounting structure to generate familiar financial statements (income statement, balance sheet, statement of cash flows, and statement of stockholders' equity), while (b) keeping the state/action space small enough to be played repeatedly by LLM agents with limited context windows.

The game runs at a \textbf{quarterly} frequency. All firms start from a \textbf{zero book state} (startup-style) and must raise capital to operate.

Crucially, the environment does \emph{not} hard-code \emph{strategic} choices by firms or the financial sector (e.g., no fixed equity multiples, no fixed revolver limits, no fixed credit-spread rules). The financial market LLM sets equity prices and credit terms directly from public information, and the firm LLM chooses operating and financing actions.

The product market is modeled by a compact preset demand system (e.g., logit shares with market-size and taste shocks), which is allowed to be hard-coded. The environment includes a \emph{bankruptcy/entry} mechanism: when a firm defaults, its assets are auctioned, proceeds pay creditors with any residual to equity, and then a fresh entrant replaces the bankrupt firm so the industry size remains constant.

Finally, the simulator must export (i) a Compustat-style \textbf{firm-quarter} panel dataset, (ii) a run-level \textbf{debrief} dataset scoring firm investment outcomes, credit performance, and equity-price efficiency, and (iii) diagnostics to flag degenerate simulations.
\end{abstract}

\section{Overview}

There are $I$ firms, indexed by $i\in\{1,\ldots,I\}$, operating for periods $t=1,\ldots,T$. 
\textbf{Each period is one fiscal quarter.} Define the fiscal calendar as
\[
\text{fyear}(t)=1+\left\lfloor\frac{t-1}{4}\right\rfloor,\qquad 
\text{fqtr}(t)=1+\big((t-1)\bmod 4\big).
\]
All \emph{flow} variables (revenues, expenses, cash flows) are quarterly; all \emph{stock} variables (balance-sheet accounts) are end-of-quarter.

\paragraph{Zero initial conditions (startup-style).}
At $t=1$, each firm begins at \emph{zero book state}:
\[
\text{Cash}=\text{AR}=\text{Inv}=\text{PPE}=\text{AP}=\text{Accr}=\text{TaxPay}=\text{RevBal}=\text{LTD}=\text{APIC}=\text{RE}=\text{TS}=0,
\]
and internal capability/brand/capacity stocks start at $A_{i0}=B_{i0}=K_{i0}=0$.
To avoid degenerate ``zero-share'' accounting, each firm has a positive share count $\text{Shares}_{i0}>0$ (e.g., $1{,}000{,}000$) with \emph{no-par stock} so the book \emph{Common stock} line item can remain $0$ at initialization.
Firms must raise capital (equity and/or debt) to fund operating and investment plans.

\paragraph{Mandatory initial capitalization (IPO phase).}
\textbf{Before} any firm can execute operating or investment decisions in its first quarter (whether at $t=1$ or upon fresh entry after bankruptcy), the environment must run a dedicated \textbf{capitalization sub-step}:
\begin{enumerate}[leftmargin=*]
    \item The firm LLM states its desired equity raise (number of new shares $\times$ offer price) and/or debt request.
    \item The financial-market LLM decides how much equity to purchase (setting the IPO price and quantity) and what credit terms to extend.
    \item The environment books the equity issuance (credit Cash, credit APIC) and any debt issuance \emph{before} the firm's operating actions are evaluated.
\end{enumerate}
If the combined capital raised is zero (financial market refuses all financing), the firm enters ``failed IPO'' status: it does not operate, records zero on all flow accounts, and is replaced by a new entrant next quarter (incrementing the incarnation index). This prevents the degenerate cycle of zero-capital firms defaulting immediately every quarter.

\paragraph{Non-negative asset invariant.}
The environment must enforce at all times, for every non-defaulted firm:
\[
\text{TotalAssets}_{it} \ge 0, \qquad \text{Cash}_{it}\ge 0.
\]
Negative total assets are \textbf{impossible} (they indicate a bookkeeping bug).  If the settlement step would cause $\text{Cash}_{it}<0$ after exhausting the revolver, the firm defaults---but the balance sheet snapshot \emph{at the moment of default} still records the cash shortfall for diagnostic purposes. No further quarterly rows are generated for a defaulted incarnation after the default quarter.

\paragraph{Where decisions live.}
\begin{itemize}[leftmargin=*]
    \item \textbf{Firms choose (LLM):} $p_{it}$ (price), $q^{\text{prod}}_{it}$ (production), $x^K_{it}$ (capex), $x^R_{it}$ (R\&D), $x^M_{it}$ (marketing), payouts, and financing requests.
    \item \textbf{Financial market chooses (LLM):} equity price $P_{it}$ and credit terms (revolver commitment, interest rates, optional rationing/clearing).
    \item \textbf{Product market (environment):} a \emph{preset} demand system maps firm prices/attributes into quantities demanded. Random macro/taste shocks are allowed and encouraged; the mapping itself is hard-coded.
    \item \textbf{Environment (non-strategic):} feasibility, settlement ordering, accounting, bankruptcy auction, and re-entry.
\end{itemize}

\paragraph{Identity discipline.}
At each quarter $t$, the simulator must output a 10-K-style:
(i) income statement, (ii) balance sheet, (iii) statement of cash flows (indirect method), and (iv) statement of stockholders' equity,
with exact reconciliation (up to numerical tolerance):
\[
\text{Assets}=\text{Liabilities}+\text{Equity},\qquad \Delta \text{Cash}=\text{CFO}+\text{CFI}+\text{CFF}.
\]
\textbf{Quarterly convention:} interest rates are treated as annualized APRs unless explicitly stated otherwise, and are converted to quarterly accruals by dividing by $4$.

\section{Agents, Information, and Objectives}

\subsection{Firm agents (LLMs)}
Each firm $i$ is controlled by an LLM that observes:
\begin{itemize}[leftmargin=*]
    \item its own private state (detailed balance sheet accounts, internal capability/brand stocks, last-period realized demand and margins);
    \item public information about rivals (prices, public quality/brand indices, and past public statements);
    \item public macro state (market size index, risk-free rate, inflation, etc.).
\end{itemize}

The firm LLM chooses actions to maximize long-run value (not hard-coded in the environment), subject to feasibility constraints.

\subsection{Financial market agent (LLM)}
A financial-market LLM observes \emph{public} information only:
\begin{itemize}[leftmargin=*]
    \item public firm financial statements (or a compressed public summary);
    \item public actions (e.g., prices, production, capex, major payouts/issuance plans);
    \item macro state (risk-free rate, volatility proxy, market size growth, etc.);
    \item past default/auction history (publicly observable outcomes).
\end{itemize}

It outputs equity prices and credit terms. The environment does not impose valuation formulas or credit limits beyond feasibility bookkeeping and numerical guardrails.

\subsection{Optional product-market agent (LLM)}
The baseline product market is a preset demand function. Optionally, a product-market LLM may supply \emph{macro/taste} shocks (e.g., total market size $M_t$ or taste shifters), but it does not ``allocate demand'' directly unless explicitly enabled. This option exists mainly to support narrative variation; the simulator must remain well-defined without it.

\section{Timing Within Each Period}

A minimal timing that supports realistic statements:

\begin{enumerate}[leftmargin=*]
    \item \textbf{Information release:} publish public summaries (prior-period statements, public indices). Provide each firm its private report.
    \item \textbf{Capitalization phase (new entrants only):} any firm in its first quarter of an incarnation (including $t=1$) runs the mandatory IPO/capitalization sub-step (see ``Mandatory initial capitalization'' in Section ``Overview''). If capitalization fails, the firm does not proceed to step 3 and is recorded as a failed IPO.
    \item \textbf{Firm actions (LLM):} each firm outputs $p_{it}$, $q^{\text{prod}}_{it}$, $x^K_{it}$, $x^R_{it}$, $x^M_{it}$, dividends, buybacks, and financing requests.
    \item \textbf{Product market clearing (preset):} demand system computes $q^{\text{dem}}_{it}$ and realized sales $q^{\text{sale}}_{it}\le q^{\text{dem}}_{it}$ subject to supply/inventory constraints.
    \item \textbf{Financial market clearing (LLM):} financial-market LLM sets equity price $P_{it}$ and credit terms for period $t$ (revolver commitment, interest rates, and any rationing/clearing of requested issuance).
    \item \textbf{Settlement and bookkeeping (with feasibility clamping):} realize revenues/costs, update working capital, compute taxes, settle financing/payouts subject to cash feasibility, draw revolver if needed, and post accounting entries.
    
    \textbf{Critical:} The environment must \emph{clamp} all LLM-requested expenditures to feasible levels before posting them.  The clamping order is:
    \begin{enumerate}[leftmargin=*]
        \item Compute available resources: beginning cash $+$ revenue collections $+$ net financing proceeds (equity issued $+$ debt issued $-$ debt repaid) $+$ available revolver commitment.
        \item Allocate to obligations in priority: (a) COGS (cost of goods actually sold), (b) interest on existing debt, (c) taxes due, (d) capex, R\&D, marketing (pro-rata if insufficient), (e) dividends and buybacks (only from remaining surplus).
        \item If total obligations exceed available resources even after full revolver draw, \emph{scale down} discretionary spending (capex, R\&D, marketing, payouts) to fit.  Record the \emph{actual} (clamped) amounts on the financial statements, not the LLM's raw requests.
    \end{enumerate}
    After clamping, $\text{Cash}_{it}\ge 0$ must hold for any non-defaulted firm.  If it still does not (e.g., mandatory costs exceed all resources), the firm defaults.
    \item \textbf{Default check and bankruptcy/entry:} if a firm is insolvent (cash infeasible after exhausting market-supplied liquidity), run an auction, distribute proceeds to claimholders, mark the firm exited, and schedule a fresh entrant for $t+1$.
    \item \textbf{Reporting:} produce 10-K-style statements and export datasets.
\end{enumerate}

\section{Product Market (Preset Demand System)}

\subsection{Public product attributes}
Each firm sells one differentiated product. The environment maintains internal ``stocks'' that proxy for capability and brand and exposes simple public indices:
quality index $Q_{it}$ and brand index $S_{it}$.

These indices are derived from internal states (e.g., accumulated R\&D and marketing with depreciation), but the mapping is fixed and transparent.

\subsection{Market size}
Let $M_t\ge 0$ be total category demand in units. In the baseline, $M_t$ follows an exogenous macro process, e.g.
\[
\log M_t = \mu_M + \rho_M \log M_{t-1} + \sigma_M \varepsilon_t,
\]
with $\varepsilon_t\sim \mathcal{N}(0,1)$ or a bounded alternative. Optionally, a product-market LLM may output $M_t$ (subject to guardrails), but the demand allocation remains preset.

\subsection{Demand allocation: multinomial logit example}
A compact baseline is multinomial logit shares:
\begin{align}
u_{it} &= \beta_Q Q_{it} + \beta_S S_{it} - \beta_P p_{it} + \xi_{it},\\
s_{it} &= \frac{\exp(u_{it})}{\sum_{j=1}^I \exp(u_{jt})},\\
q^{\text{dem}}_{it} &= s_{it}\, M_t,
\end{align}
where $\xi_{it}$ are taste shocks (drawn by the environment or optionally suggested by a product-market LLM). This section is illustrative; any small, documented preset demand function is acceptable.

\subsection{Realized sales under supply constraints}
Let beginning inventory be $\text{Inv}_{i,t-1}$ and production $q^{\text{prod}}_{it}$. Then feasible sales satisfy
\[
q^{\text{sale}}_{it}\le q^{\text{dem}}_{it},\qquad q^{\text{sale}}_{it}\le \text{Inv}_{i,t-1}+q^{\text{prod}}_{it}.
\]
Unsatisfied demand is lost (baseline) and excess supply becomes inventory.

\section{Operations Technology and State Dynamics}

\subsection{Production and costs}
Let unit variable cost depend on capability:
\[
c_{it} = c_0 \exp(-\phi A_{it}),
\]
where $A_{it}$ is a capability stock (internal, not directly observed by rivals). Then
\[
\text{COGS}_{it}=c_{it}\, q^{\text{sale}}_{it}.
\]
Fixed operating cost $F$ enters SG\&A.

\subsection{Capability and brand accumulation}
Internal stocks evolve as
\begin{align}
A_{i,t} &= (1-\delta_A)A_{i,t-1} + \eta_A x^R_{it},\\
B_{i,t} &= (1-\delta_B)B_{i,t-1} + \eta_B x^M_{it},
\end{align}
and public indices are monotone transforms, e.g.
\[
Q_{it}=\log(1+A_{it}),\qquad S_{it}=\log(1+B_{it}).
\]

\subsection{Capacity and PPE}
Capacity $K_{it}$ evolves with capex net of depreciation:
\[
K_{it} = (1-\delta_K)K_{i,t-1} + \eta_K x^K_{it}.
\]
Book PPE is tracked as $\text{PPE}_{it}$ with depreciation $\text{Dep}_{it}=\delta_P \text{PPE}_{i,t-1}$.

\section{Working Capital, Taxes, and Cash Mechanics}

\subsection{Working capital structure (quarterly)}
Use a compact working-capital model at the quarterly frequency:
\begin{itemize}[leftmargin=*]
    \item Accounts receivable: $\text{AR}_{it} = \theta_{AR}\, \text{NetSales}_{it}$ (one-quarter collection lag).
    \item Inventory: $\text{Inv}_{it} = \text{Inv}_{i,t-1} + q^{\text{prod}}_{it} - q^{\text{sale}}_{it}$ (in units) times a unit-cost valuation.
    \item Accounts payable: $\text{AP}_{it} = \theta_{AP}\, \text{COGS}_{it}$ (one-quarter payment lag).
    \item Accrued expenses: $\text{Accr}_{it}=\theta_{Accr}(\text{SG\&A}_{it}+\text{R\&D}_{it})$ (lagged).
\end{itemize}

\subsection{Taxes (quarterly accrual, lagged payment)}
Corporate tax expense (quarterly):
\[
\text{TaxExpense}_{it}=\tau \max(\text{PreTaxIncome}_{it},0).
\]
Taxes payable can be lagged:
\[
\text{TaxPay}_{it} = \text{TaxPay}_{i,t-1} + \text{TaxExpense}_{it} - \text{TaxPaid}_{it},
\]
with $\text{TaxPaid}_{it}$ set by a simple policy (e.g., pay last quarter's taxes, subject to cash feasibility).

\subsection{Cash update (conceptual)}
Cash updates through operating, investing, and financing cash flows as in the statement of cash flows (Section~\ref{sec:scf}).

\section{Financing and Financial Market (LLM-Determined)}

\subsection{Financial-market LLM outputs (quarterly contracting)}
For each firm, the financial-market LLM outputs (publicly) at least:
\begin{itemize}[leftmargin=*]
    \item Equity price $P_{it}>0$ (interpreted as the \emph{start-of-quarter} price for quarter $t$).
    \item Revolver terms: commitment $L^{\text{commit}}_{it}\ge 0$ and annualized rate $r^{\text{rev}}_{it}\ge 0$.
    \item Term debt terms: annualized rate $r^{\text{term}}_{it}\ge 0$ and optional max-new amount $L^{\text{term}}_{it}\ge 0$.
    \item Optional clearing: how much of requested equity/debt issuance and buybacks are filled.
\end{itemize}
The environment validates only numerical sanity and firm-ID consistency (no extra/missing firms).

\subsection{Interest and debt accounting (quarterly accrual)}
Track short-term borrowings (revolver balance) $\text{RevBal}_{it}$ and long-term debt $\text{LTD}_{it}$.
Treat LLM-provided rates as \emph{annualized APRs} and accrue interest quarterly using a $1/4$ factor:
\[
\text{IntExp}_{it} = \frac{r^{\text{rev}}_{it}}{4}\,\text{RevBal}_{i,t-1} + \frac{r^{\text{term}}_{it}}{4}\,\text{LTD}_{i,t-1}.
\]
(Using beginning-of-quarter balances is an acceptable approximation; the code may use average balances if desired.)
Debt issuance/repayment changes balances per clearing.

\subsection{Equity}
Shares outstanding $\text{Shares}_{it}$ change with issuance/buybacks. Equity market capitalization is
\[
\text{MktCap}_{it} = P_{it}\,\text{Shares}_{it}.
\]
The environment does not impose any mapping from earnings/book value to $P_{it}$.

\section{Accounting System and Financial Statements (10-K Style)}

The environment is a bookkeeping engine that produces a familiar, 10-K-like presentation using a compact chart of accounts.
The goal is not to reproduce all footnote complexity, but to ensure that (i) labels look realistic and (ii) the core identities reconcile.

\subsection{Income statement (statement of operations)}
Define
\begin{align}
\text{GrossProfit}_{it}&=\text{NetSales}_{it}-\text{COGS}_{it},\\
\text{R\&D}_{it}&=x^R_{it},\\
\text{SG\&A}_{it}&=F+x^M_{it},\\
\text{EBIT}_{it}&=\text{GrossProfit}_{it}-\text{R\&D}_{it}-\text{SG\&A}_{it}-\text{Dep}_{it},\\
\text{OtherIncome}_{it}&=0\ \text{ (baseline)}.
\end{align}
Net income is
\begin{equation}
\text{NetIncome}_{it}=\text{PreTaxIncome}_{it}-\text{TaxExpense}_{it}.
\end{equation}
A 10-K-like presentation:
\begin{center}
\begin{tabular}{l r}
\toprule
Net sales (revenue) & $\text{NetSales}_{it}$\\
Cost of revenue & $-\text{COGS}_{it}$\\
\midrule
Gross profit & $\text{GrossProfit}_{it}$\\
Operating expenses: & \\
\quad Research and development & $-\text{R\&D}_{it}$\\
\quad Selling, general and administrative & $-\text{SG\&A}_{it}$\\
\quad Depreciation and amortization & $-\text{Dep}_{it}$\\
\midrule
Operating income (loss) & $\text{EBIT}_{it}$\\
Interest expense & $-\text{IntExp}_{it}$\\
Other income (expense), net & $\text{OtherIncome}_{it}$\\
\midrule
Income (loss) before income taxes & $\text{PreTaxIncome}_{it}$\\
Provision for income taxes & $-\text{TaxExpense}_{it}$\\
\midrule
Net income (loss) & $\text{NetIncome}_{it}$\\
\bottomrule
\end{tabular}
\end{center}

\subsection{Balance sheet}
Key asset accounts:
\begin{align}
\text{Cash}_{it} &\ge 0 \text{ (strictly enforced; violation triggers default)},\\
\text{AR}_{it}&=\theta_{AR}\,\text{NetSales}_{it},\\
\text{Inv}_{it}&\ge 0,\\
\text{PPE}_{it}&=(1-\delta_P)\text{PPE}_{i,t-1}+x^K_{it} \ge 0.
\end{align}
\textbf{Important:} since all asset accounts are non-negative, $\text{TotalAssets}_{it}\ge 0$ follows automatically. If the accounting engine ever produces $\text{TotalAssets}_{it}<0$, this is a \emph{bug} in settlement/clamping logic, not a valid economic outcome.
Key liability accounts:
\begin{align}
\text{AP}_{it}&=\theta_{AP}\,\text{COGS}_{it},\\
\text{Accr}_{it}&=\theta_{Accr}(\text{SG\&A}_{it}+\text{R\&D}_{it}),\\
\text{TaxPay}_{it}&\ge 0,\\
\text{RevBal}_{it}&\ge 0,\qquad \text{LTD}_{it}\ge 0.
\end{align}

Stockholders' equity (book) is the residual:
\[
\text{EquityBook}_{it}=\text{TotalAssets}_{it}-\text{TotalLiab}_{it}.
\]
For reporting, represent equity with compact components (common stock, APIC, treasury stock, retained earnings, AOCI baseline 0) that reconcile to the residual.

A 10-K-like balance sheet presentation:
\begin{center}
\begin{tabular}{l r}
\toprule
\textbf{Assets} & \\
Cash and cash equivalents & $\text{Cash}_{it}$\\
Accounts receivable, net & $\text{AR}_{it}$\\
Inventory & $\text{Inv}_{it}$\\
Prepaid expenses and other current assets & $\text{Prepaid}_{it}$\\
\midrule
Total current assets & $\text{CA}_{it}$\\
Property and equipment, net & $\text{PPE}_{it}$\\
Other assets & $\text{OtherA}_{it}$\\
\midrule
Total assets & $\text{TotalAssets}_{it}$\\
\midrule
\textbf{Liabilities and stockholders' equity} & \\
Accounts payable & $\text{AP}_{it}$\\
Accrued expenses and other current liabilities & $\text{Accr}_{it}$\\
Income taxes payable & $\text{TaxPay}_{it}$\\
Short-term borrowings (revolver) & $\text{RevBal}_{it}$\\
\midrule
Total current liabilities & $\text{CL}_{it}$\\
Long-term debt & $\text{LTD}_{it}$\\
Other long-term liabilities & $\text{OtherL}_{it}$\\
\midrule
Total liabilities & $\text{TotalLiab}_{it}$\\
\midrule
Common stock & $\text{CS}_{it}$\\
Additional paid-in capital & $\text{APIC}_{it}$\\
Treasury stock & $-\text{TS}_{it}$\\
Retained earnings (accumulated deficit) & $\text{RE}_{it}$\\
Accumulated other comprehensive income (loss) & $\text{AOCI}_{it}$\\
\midrule
Total stockholders' equity & $\text{TotalEquity}_{it}$\\
\midrule
Total liabilities and stockholders' equity & $\text{TotalLiab}_{it}+\text{TotalEquity}_{it}$\\
\bottomrule
\end{tabular}
\end{center}

\subsection{Statement of cash flows (indirect method)}\label{sec:scf}
Cash flow from operations:
\begin{align}
\text{CFO}_{it} &= \text{NetIncome}_{it} + \text{Dep}_{it} \\
&\quad -\Delta \text{AR}_{it} - \Delta \text{Inv}_{it} + \Delta \text{AP}_{it} + \Delta \text{Accr}_{it} + \Delta \text{TaxPay}_{it}.
\end{align}
Cash flow from investing:
\[
\text{CFI}_{it} = -x^K_{it}.
\]
Cash flow from financing:
\begin{align}
\text{CFF}_{it} &= +\text{EquityIssuance}_{it} - \text{Buyback}_{it} - \text{Div}_{it} \\
&\quad + \Delta \text{LTD}_{it} + \Delta \text{RevBal}_{it}.
\end{align}
Cash reconciliation:
\[
\text{Cash}_{it}=\text{Cash}_{i,t-1}+\text{CFO}_{it}+\text{CFI}_{it}+\text{CFF}_{it}.
\]
10-K-like presentation should list the above line items.

\subsection{Statement of stockholders' equity}
Retained earnings roll forward:
\[
\text{RE}_{it}=\text{RE}_{i,t-1}+\text{NetIncome}_{it}-\text{Div}_{it}.
\]
Treasury stock under cost method increases with buybacks and decreases with re-issuance (if allowed).
AOCI baseline is 0 unless additional features are added.

\section{Feasibility, Default, Bankruptcy Auction, and Entry}

\subsection{Bounds and action constraints (game-friendly)}
To keep the game playable, enforce simple \emph{numerical/physical} bounds:
\begin{itemize}[leftmargin=*]
    \item Production: $0\le q^{\text{prod}}_{it}\le \kappa K_{i,t-1}$.
    \item Spending: $x^K_{it},x^R_{it},x^M_{it}\ge 0$ (optionally capped at a large technical limit).
    \item Payout requests: $\text{Div}_{it},\text{Buyback}_{it}\ge 0$.
    \item Financing requests: $\text{DebtIssue}_{it},\text{DebtRepay}_{it},\text{EquityIssuance}_{it}\ge 0$.
\end{itemize}

\subsection{Cash feasibility and revolver (market-supplied commitment)}
At settlement, if projected cash would fall below zero, the environment draws the revolver \emph{up to the available commitment supplied by the financial-market LLM}.
If the financial-market LLM supplies a per-period commitment $L^{\text{commit}}_{it}$, then the environment enforces
\[
\text{RevBal}_{it}\le L^{\text{commit}}_{it}.
\]
There is no exogenous hard-coded revolver limit besides technical sanity bounds.

\subsection{Default trigger}
Default rule:
\[
\text{Cash}_{it}<0 \text{ after exhausting market-supplied liquidity} \quad \Rightarrow\quad \text{Default}.
\]

\subsection{Bankruptcy auction (liquidation) and claim recovery}
Upon default at the end of quarter $t$:
\begin{enumerate}[leftmargin=*]
    \item \textbf{Auction proceeds:} sell assets at recovery rates (salvage haircuts) applied to book values:
    \begin{align}
    \text{Proceeds}_{it} &= \lambda_C \text{Cash}_{it} + \lambda_{AR}\text{AR}_{it} + \lambda_{Inv}\text{Inv}_{it} + \lambda_{PPE}\text{PPE}_{it} + \lambda_{O}\text{OtherA}_{it}.
    \end{align}
    Baseline: $\lambda_C=1$, and $\lambda_{AR},\lambda_{Inv},\lambda_{PPE},\lambda_O\in[0,1]$.
    \item \textbf{Pay liabilities in a simple priority waterfall:}
    \begin{enumerate}[leftmargin=*]
        \item Revolver (short-term borrowings),
        \item Term debt (long-term debt),
        \item Trade and statutory payables (AP, accrued liabilities, taxes).
    \end{enumerate}
    Any unpaid balance is a loss to the respective claimant class.
    \item \textbf{Residual to equity:} if proceeds exceed total liabilities, residual value is distributed to equity holders as liquidation proceeds; otherwise equity receives zero.
\end{enumerate}
This mechanism ensures that bankruptcy outcomes affect realized creditor returns (used for scoring the financial-market LLM).

\subsection{Fresh entry after bankruptcy (zero reset)}
To keep the industry size constant, if firm $i$ defaults in $t$, then in $t+1$ a \emph{fresh entrant} replaces it:
\begin{itemize}[leftmargin=*]
    \item The entrant begins from the same \textbf{zero book state} as at $t=1$ (Section ``Overview'') and with internal stocks reset to $A=B=K=0$.
    \item The entrant inherits the \emph{slot} $i$ but receives a new \emph{incarnation index} $g$ so the internal key is $(i,g)$.
    \item The entrant also receives a new \emph{agent fingerprint} (Section ``LLM Interface''): a random but fixed ``style''/risk profile and nonce tokens, so entrants are not behavioral clones.
    \item Publicly, the identifier can remain ``Firm $i$'' while internal files and datasets use $(i,g)$ to prevent hallucinated extra firms.
\end{itemize}

\paragraph{Preventing death spirals (consecutive-failure cap).}
If a firm slot $i$ has $N_{\max}^{\text{fail}}$ (e.g., 3) consecutive incarnations that default in their first quarter (``failed IPO'' or immediate default), the environment should:
\begin{enumerate}[leftmargin=*]
    \item Pause the slot for one quarter (the slot is vacant, industry has $I-1$ active firms),
    \item On the next entry attempt, provide the financial-market LLM with an explicit flag noting the slot's history of failed capitalization,
    \item If the pattern persists beyond $2 N_{\max}^{\text{fail}}$ total consecutive failures, freeze the slot permanently for the remainder of the run and log a diagnostic warning.
\end{enumerate}
This prevents the degenerate outcome of an incarnation counter climbing to $30+$ with zero activity each quarter.

\section{Data Outputs: Compustat-Style Quarterly Panel, Debrief, and Diagnostics}

\subsection{Compustat-style firm-quarter dataset (Compustat Quarterly flavor)}
The simulator must export a CSV panel with one row per \emph{firm-incarnation per quarter}:
\[
(\text{run\_id},\ i,\ g,\ t)\ \mapsto\ \text{Compustat-style quarterly variables},
\]
where $t$ is a fiscal quarter and $(\text{fyear},\text{fqtr})$ is defined in Section ``Overview''.

Include at least the following columns (approximate Compustat naming; suffix ``q'' denotes quarterly):
\begin{itemize}[leftmargin=*]
    \item Identifiers: \texttt{run\_id}, \texttt{firm\_id}, \texttt{incarnation}, \texttt{fyearq}, \texttt{fqtr}, \texttt{datadate}.
    \item Balance sheet (end-of-quarter): 
    \texttt{atq} (total assets), \texttt{ltq} (total liabilities), \texttt{ceqq} (common equity), 
    \texttt{cheq} (cash), \texttt{rectq} (AR), \texttt{invtq} (inventory), \texttt{ppentq} (PPE, net), 
    \texttt{apq} (accounts payable), \texttt{acoq} (accrued expenses and other current liabilities),
    \texttt{lctq} (current liabilities), \texttt{dlcq} (short-term debt / revolver), 
    \texttt{dlttq} (long-term debt), \texttt{txpq} (taxes payable).
    
    \textbf{Component reconciliation:} $\texttt{lctq} = \texttt{apq} + \texttt{acoq} + \texttt{txpq} + \texttt{dlcq}$, and $\texttt{atq} = \texttt{cheq} + \texttt{rectq} + \texttt{invtq} + \texttt{ppentq} + \text{(other assets)}$. These must hold exactly (up to rounding).
    \item Income statement (quarterly flows): 
    \texttt{saleq} (net sales), \texttt{cogsq}, \texttt{gpq} (gross profit),
    \texttt{xrdq}, \texttt{xsgaq}, \texttt{dpq}, \texttt{oiadpq} (operating income after dep),
    \texttt{xintq}, \texttt{piq} (pretax income), \texttt{txtq} (tax expense), \texttt{niq}.
    \item Cash flow (quarterly): 
    \texttt{oancfq} (operating cash flow), \texttt{capxq}, 
    \texttt{dltisq} (debt issued), \texttt{dltrq} (debt repaid),
    \texttt{sstkq} (stock issued), \texttt{prstkcq} (stock repurchased),
    \texttt{dvq} (dividends), \texttt{chechq} (change in cash).
    \item Market: \texttt{prccq} (equity price at start of quarter), \texttt{cshoq} (shares), \texttt{mkvaltq} (market cap).
    \item Events and pricing diagnostics: \texttt{default\_flag}, \texttt{auction\_proceeds}, 
    \texttt{debt\_recovery\_rate}, \texttt{equity\_recovery}, and \texttt{pricing\_error\_q} (defined below).
\end{itemize}
The mapping is approximate but labels should match Compustat conventions and be consistent across runs.

\paragraph{Panel uniqueness and validation.}
The key $(\texttt{run\_id},\texttt{firm\_id},\texttt{incarnation},\texttt{fyearq},\texttt{fqtr})$ must be \textbf{unique} in the panel.  The export code must deduplicate before writing.  Additionally, the export must validate the following invariants for every non-default row (reject or warn if violated):
\begin{itemize}[leftmargin=*]
    \item $\texttt{atq}\ge 0$ (total assets non-negative),
    \item $\texttt{cheq}\ge 0$ (cash non-negative for non-defaulted firms),
    \item $|\texttt{atq}-\texttt{ltq}-\texttt{ceqq}|<1$ (balance-sheet identity),
    \item $|\texttt{lctq}-(\texttt{apq}+\texttt{acoq}+\texttt{txpq}+\texttt{dlcq})|<1$ (current liabilities reconcile with components).
\end{itemize}
For the default quarter's row, negative cash is permitted as a diagnostic record, but \texttt{atq} should still be computed correctly as the sum of all asset accounts.

\paragraph{Pricing error per firm-quarter.}
Let $P_{it}$ denote the \emph{start-of-quarter} price set by the financial-market LLM for quarter $t$, and let $\text{Div}_{it}$ be the dividend paid \emph{during} quarter $t$ (end-of-quarter payout). Let $r^f_t$ be the annualized risk-free rate; define the quarterly discount factor as $1+r^f_t/4$.
Define an ex post one-quarter ``fundamental'' (realized discounted payoff) as
\[
P^{\star}_{it}=\frac{\text{Div}_{it}+P_{i,t+1}}{1+r^f_t/4},
\]
and record the pricing error
\[
\text{pricing\_error\_q}_{it}=P_{it}-P^{\star}_{it},
\]
for all $t<T$ (for $t=T$ omit or use a terminal convention).

\subsection{Debrief dataset (auto-expanding across simulations)}
Each run must produce a \textbf{debrief CSV} that appends across runs (create file if absent, append if present). 
The debrief file is the main ``scoreboard'' for repeated simulations and must include:

\paragraph{Firm investment performance.}
For each firm-incarnation $(i,g)$, report:
\begin{itemize}[leftmargin=*]
    \item Terminal total claim value:
    \[
    V^{\text{end}}_{i,g}=
    \begin{cases}
    P_{iT}\,\text{Shares}_{iT} + \text{RevBal}_{iT}+\text{LTD}_{iT} & \text{if alive at }T,\\
    \text{Proceeds}_{i,g} & \text{if defaulted before }T.
    \end{cases}
    \]
    \item Equity-investor cash-flow series (quarterly): 
    \[
    \text{CF}^{E}_{it}= \text{Div}_{it}+\text{BuybackPaid}_{it}-\text{EquityIssueProceeds}_{it}
    \]
    plus a terminal value at exit:
    $+\ P_{iT}\text{Shares}_{iT}$ if alive at $T$, or $+\ \text{EquityResidual}_{i,g}$ if liquidated.
    From these cash flows compute \textbf{equity IRR} and \textbf{equity total return}.
    \item Simple summary components: total capital raised (equity + debt), total distributions to equity (dividends + buybacks + liquidation residual), default quarter (if any), lifespan in quarters.
\end{itemize}

\paragraph{Credit performance (financial intermediaries).}
Treat the financial-market agent as a lender supplying revolver draws and term debt. Record quarterly lender cash flows:
\begin{itemize}[leftmargin=*]
    \item Principal outflows when credit is extended (revolver draws and term debt issuance),
    \item Principal inflows when repaid,
    \item Interest inflows (quarterly accrual based on annualized rates divided by 4),
    \item Bankruptcy recoveries from the liquidation waterfall.
\end{itemize}
Compute and report (for the whole credit portfolio and optionally by instrument):
\textbf{total return}, \textbf{IRR}, \textbf{loss rate}, and \textbf{default rate}.

\paragraph{Equity price efficiency (financial-market agent).}
Aggregate pricing errors $\text{pricing\_error\_q}_{it}$ into run-level metrics such as:
\begin{itemize}[leftmargin=*]
    \item MAPE or sMAPE of $P_{it}$ vs. $P^{\star}_{it}$,
    \item RMSE and bias of pricing errors,
    \item correlation between $P_{it}$ and $P^{\star}_{it}$.
\end{itemize}
These constitute the financial-market LLM's ``price efficiency'' score.

\subsection{Diagnostics: non-degenerate behavior checks}
The simulator should emit a simple diagnostics report (CSV/JSON/Markdown) to help users judge realism. At minimum, warn/flag runs where:
\begin{itemize}[leftmargin=*]
    \item all firms choose identical actions for many quarters (prices, capex, R\&D, marketing nearly identical);
    \item investment is persistently zero (e.g., $\text{capxq}=0$ and $\text{xrdq}=0$ for all firms over long spans);
    \item market shares are constant with no entry/exit despite shocks (unless shocks are disabled).
\end{itemize}
Provide guidance to inspect the Compustat-style panel for dispersion in \texttt{saleq}, \texttt{capxq}, \texttt{xrdq}, margins, leverage, and stock returns.

\paragraph{Hard invariant checks (must pass or the run is flagged invalid).}
The diagnostics must also verify:
\begin{enumerate}[leftmargin=*]
    \item \textbf{No negative total assets} for non-default rows: $\texttt{atq}\ge 0$.
    \item \textbf{No negative cash} for non-default rows: $\texttt{cheq}\ge 0$.
    \item \textbf{Balance-sheet identity}: $|\texttt{atq}-\texttt{ltq}-\texttt{ceqq}|<1$ for every row.
    \item \textbf{Cash reconciliation}: $|\texttt{chechq}-(\text{Cash}_t - \text{Cash}_{t-1})|<1$ for consecutive quarters of the same incarnation.
    \item \textbf{Current-liabilities decomposition}: $|\texttt{lctq}-(\texttt{apq}+\texttt{acoq}+\texttt{txpq}+\texttt{dlcq})|<1$.
    \item \textbf{No duplicate keys}: each $(\texttt{run\_id},\texttt{firm\_id},\texttt{incarnation},\texttt{fyearq},\texttt{fqtr})$ appears exactly once.
    \item \textbf{Death-spiral detection}: flag if any firm slot reaches incarnation $>N_{\max}^{\text{fail}}$ consecutive first-quarter defaults.
\end{enumerate}

\section{LLM Interface, Prompting, Memory, Randomness, and Hallucination Guardrails}

\subsection{Structured outputs and firm-ID guardrails}
All LLMs must output \emph{JSON only}. Prompts must include the authoritative list of firms (e.g., \texttt{["firm\_1",...,"firm\_I"]}) and a schema requiring keys only from that set. The environment must validate:
\begin{itemize}[leftmargin=*]
    \item no missing firms in required fields;
    \item no extra firm keys (reject hallucinated firms);
    \item numeric sanity (finite; nonnegative where required).
\end{itemize}
Invalid outputs trigger a bounded repair loop (retry with schema and the offending output).

\subsection{Within-run memory system (mandatory)}
Because this is a multi-quarter game, prompts must include \emph{what happened previously}. The environment maintains two memories:

\paragraph{Market memory (public, canonical).}
A structured event log for each quarter including, at minimum:
market size $M_t$, macro shocks, each firm's public actions (prices, payouts, issuance), realized sales/market shares, credit terms, equity prices/returns, and any defaults/auctions/entries.
This is the canonical record used to prevent hallucinated events.

\paragraph{Firm memory (private, canonical).}
For each firm-incarnation $(i,g)$, maintain a private record of its own internal states and outcomes (capability/brand/capacity stocks, unit costs, inventory, and any internal KPIs).

\subsection{Temporal prompts with memory management}
Each agent prompt (firms and financial market) should include:
\begin{itemize}[leftmargin=*]
    \item \textbf{Current quarter header:} $(\text{fyear},\text{fqtr})$, macro state, and the authoritative firm list.
    \item \textbf{Short memory:} last $K$ quarters of structured market-memory events (and, for firms, their private events).
    \item \textbf{Long memory:} a compressed summary of earlier quarters, updated each quarter \emph{only from the canonical log} (deterministic summarization preferred; LLM summarization allowed only if constrained to summarize provided facts).
\end{itemize}
When prompt length approaches a limit, drop older raw quarters but keep the long-memory summary.

\subsection{Randomness, agent fingerprints, and reproducibility}
To avoid degenerate ``all firms behave the same'' outcomes, the environment includes exogenous randomness and agent heterogeneity:
\begin{itemize}[leftmargin=*]
    \item \textbf{Exogenous shocks:} macro shocks (e.g., market size, risk-free rate) and firm taste shocks $\xi_{it}$ in demand (Section ``Product Market'').
    \item \textbf{Agent fingerprints (fixed within a run):} each firm agent and the financial-market agent receive a random but fixed ``profile'' (e.g., risk appetite, growth-vs-profitability tilt) and a nonce token. This profile is included in the system prompt so that different agents make different choices even under similar states.
    \item \textbf{Determinism option:} for reproducibility, use temperature $=0$ and (if supported) pass an explicit RNG seed to the local LLM runtime. Randomness then enters only through the environment's seeded shocks and agent fingerprints.
\end{itemize}

\subsection{Consulting a database of past simulations (cross-run memory)}
To improve decisions, the simulator may provide the firm and financial-market LLMs with \emph{retrieved cases} from prior runs stored as CSVs (Compustat-style quarterly panel and debrief dataset). A simple retrieval rule:
\begin{itemize}[leftmargin=*]
    \item compute a feature vector (e.g., leverage, profitability, growth, liquidity);
    \item select the $K$ nearest historical \emph{firm-quarters};
    \item include their actions/outcomes as ``prior episodes'' in the prompt.
\end{itemize}
The code must: (i) create the historical CSV files if absent, and (ii) append new runs if present.

\section{Baseline Parameterization (Example)}

Choose parameters to produce plausible statement magnitudes while keeping the game stable at a \textbf{quarterly} frequency:
\begin{itemize}[leftmargin=*]
    \item $I=4$ firms; $T=80$ quarters (20 fiscal years) as a baseline.
    \item Shares at start: $\text{Shares}_{i0}=1{,}000{,}000$ (no-par; book common stock can be $0$).
    \item Demand parameters: $\beta_P>0$ large enough for price sensitivity; $\beta_Q,\beta_S>0$.
    \item Working capital lags: $\theta_{AR}\in[0.05,0.30]$, $\theta_{AP}\in[0.05,0.30]$, $\theta_{Accr}\in[0.05,0.30]$.
    \item Annual tax rate $\tau\approx 0.21$ applied to quarterly pretax income.
    \item Bankruptcy recovery rates: $\lambda_{AR},\lambda_{Inv},\lambda_{PPE}\in[0.2,0.9]$.
    \item Exogenous shock scales: $\sigma_M$ (market size volatility) and taste shock scale for $\xi_{it}$ chosen to generate nontrivial market-share movement.
    \item \textbf{Consecutive-failure cap}: $N_{\max}^{\text{fail}}=3$ (configurable).
\end{itemize}

\paragraph{Prompt guidance for LLM agents (spending reasonableness).}
The firm-agent system prompt should include guidance that:
\begin{itemize}[leftmargin=*]
    \item Total quarterly spending (COGS $+$ R\&D $+$ SG\&A $+$ capex) should not routinely exceed total available capital (cash $+$ credit). The environment will clamp if it does, but the LLM should aim for feasible plans.
    \item A newly capitalized startup should not spend its entire capital in one quarter. Prompt the LLM with a ``capital budget'' derived from its balance sheet (e.g., ``You have \$X in cash and \$Y in credit; plan accordingly.'').
    \item Dividends and buybacks should only be attempted when retained earnings and cash are positive.
\end{itemize}

\end{document}
